\section{Обход графа в ширину}

Обход графа в ширину (Breadth-First Search, BFS)~--- это одна из фундаментальных задач анализа графов, для решения которой существуют хорошо известные алгоритмы, изложенные в классической литературе\sidenote{
    Например, в книге <<Алгоритмы. Построение и анализ>>~\cite{кормен2009алгоритмы} Томаса Кормена и других.
}.
Более того, данный алгоритм является основой для многих алгоритмов поиска в графе.

В общих чертах, задача заключается в том, чтобы начиная с некоторой заданной вершины графа (источника) обойти все достижимые из неё вершины в некотором порядке.
Обход в ширину решает эту задачу и задаёт порядок посещения вершин с использованием \emph{фронта обхода} (или просто \emph{фронта}).
\begin{definition}[Фронт обхода]
    \emph{Фронт обхода} графа $\mbfscrG = \langle V, E, L \rangle$~--- это вектор размера $|V|$, содержащий информацию%
    \sidenote{В самом простом случае это будет булев вектор $F$ такой, что $F[i] = 1$ тогда и только тогда, когда \circled{$v_i$} достижима из стартовой вершины за заданное число шагов.}
    о вершинах графа, достижимых из стартовой ровно за $k$ шагов для заданного $k$.
    Иными словами, все вершины, информация о которых есть во фронте, достижимы из стартовой вершины за одинаковое число шагов.
\end{definition}

Один шаг алгоритма заключается в рассмотрении всех вершин, смежных со всеми вершинами текущего фронта и формировании из них нового фронта.
При формировании нового фронта необходимо учесть, что в процессе обхода не нужно посещать одну и ту же вершину несколько раз%
\sidenote{Как правило не нужно.
    Существуют алгоритмы, использующие аналогичную конструкцию с фронтом, но посещающие некоторые вершины несколько раз.
    Например, поиск кратчайших путей от заданной вершины.}.
Схематичное изображение такого шага приведено на рисунке~\ref{fig:bfs_schema}.
\begin{marginfigure}
        \begin{tikzpicture}
        \begin{scope}
            \node [circle, draw] (p1)                             {};
            \node [text width=0.1cm] (p2) [below = 0.3 of p1]     {$\vdots$};
            \node [circle, draw] (p3) [below = 0.3 of p2]         {};
            \node [text width=0.1cm] (p4) [below = 0.3 of p3]     {$\vdots$};
            \node [text width=0.1cm] (p5) [left = 0.8 of p2]     {$\vdots$};
            \begin{pgfonlayer}{background}
                \path[fill=green,rounded corners]
                ([yshift=5pt,xshift=-2pt] p1.north west) rectangle ([yshift=-5pt,xshift=2pt] p4.south east);
            \end{pgfonlayer}
        \end{scope}

        \begin{scope}
            \node [circle, draw] (q1) [right = of p1]                          {};
            \node [text width=0.1cm] (q2)     [below = 0.3 of q1]              {$\vdots$};
            \node [text width=0.1cm] (q4)     [right = 0.8 of q2]              {$\vdots$};
            \node [circle, draw] (q3) [below = 0.3 of q2]                      {};
            \begin{pgfonlayer}{background}
                \path[fill=yellow,rounded corners]
                ([yshift=5pt,xshift=-2pt] q1.north west) rectangle ([yshift=-5pt,xshift=2pt] q3.south east);
            \end{pgfonlayer}
        \end{scope}
        \node[text width=2.5cm] (new) [below right = 0.4 of q3] {\scriptsize{новый фронт}};
        \node[text width=2.5cm] (current) [below right = 0.1 of p4] {\scriptsize{текущий фронт}};

        \path[->]
        (p1) edge (q1)
        (p1) edge (q3)
        (p3) edge (q3)
        (p2) edge (q2)
        (p4) edge (q2)
        (p5) edge (p2)
        (p5) edge (p4)
        (q2) edge (q4)
        (new) edge [dashed] ([yshift=-5pt,xshift=-2.5pt] q3.south east)
        (current) edge [dashed] ([yshift=2.5pt, xshift=0.5pt] p4.south east)
        ;
    \end{tikzpicture}

    \caption{Схема одного шага обхода в ширину}
    \label{fig:bfs_schema}
\end{marginfigure}

Алгоритм обхода в ширину может быть переформулирован в терминах матрично-векторных операций следующим образом%
\sidenote{
    Стоит отметить, что ситуация с не менее известным обходом в глубину (Depth-First Search, DFS) более сложная: на момент написания текста не известно <<естественного>> выражения данного обхода в терминах линейной алгебры.
    Доказательство невозможности такого построения также не предъявлены.
    При этом, решения для частных случаев (деревья, ориентированные графы без циклов) предложены, например, в работе~\cite{10.1145/3315454.3329962}.}.
Пусть фронт~--- вектор размера $n$, а сам граф представлен матрицей смежности.
Тогда один шаг~--- просмотр всех смежных вершин для формирования нового нового фронта~--- это умножение текущего фронта на матрицу смежности.
Для того, чтобы отслеживать уже посещённые вершины, поддерживается булев вектор \emph{visited} размера $|V|$, такой, что $\emph{visited}[i] = 1$ тогда и только тогда, когда  \circled{$v_i$} уже посещалась в процессе обхода\sidenote{
    Делать вектор \emph{visited} булевым удобно только в самых простых случаях.
    В общем случае он будет хранить некоторую информацию о вершине.
    Например, номер шага, на котором её посетили, её предка, или даже расстояние до неё от стартовой вершины.
}.
Данный вектор используется как маска для фильтрации кандидатов в новый фронт.

Такие шаги повторяются до тех пор, пока фронт не окажется пустым ($\overline{0}$, нет новых вершин для посещения).
А результатом работы является информация, накопленная в \emph{visited}.

Псевдокод самой простой версии алгоритма\sidenote{
    Возможны вариации. Точно не самый оптимальный. Ссылка на AGraph и на статью какую-нибудь!!!!
} представлен на листинге~\ref{algo:BFS_linal}.
Нам понадобятся несколько алгебраических структур.
Во-первых, $\BbbB = \langle \{0,1\},\vee, \wedge \rangle$~--- стандартное булево полукольцо.
Во-вторых, $\BbbM = \langle \{0,1\},\oplus \rangle$, где $\oplus$ определена следующим образом\sidenote{
    Инвертированная маска: в результат попадают только те значения из первого вектора, которым соответствуют нулевые значения во втором.
}:
\begin{minipage}{0.25\textwidth}
    \begin{itemize}
        \item $0 \oplus 0 = 0$;
        \item $1 \oplus 1 = 0$;
    \end{itemize}
\end{minipage}
\begin{minipage}{0.25\textwidth}
    \begin{itemize}
        \item $0 \oplus 1 = 0$;
        \item $1 \oplus 0 = 1$.
    \end{itemize}
\end{minipage}

\mytodo{Вынести обозначения операций в алгебраических струткруах (типа $\mmult{\BbbB}$)  в соответствующие разделы линейной алгебры. Перенести теховские определения в правильный файл из style}

\begin{algorithm}
    %\SetAlgoLined
    \KwData{$M$~--- булева матрица смежности графа\;
        $k$~--- номер стартовой вершины}
    \KwResult{Вектор, указывающий, какие вершины достижимы из стартовой}
    $\emph{current\_front} \leftarrow \overline{0}$\;
    $\emph{visited} \leftarrow \overline{0}$\;
    $\emph{current\_front}[k] \leftarrow 1$\;
    \While{$\emph{current\_front} \neq \ \overline{0}$}{
        $\emph{visited} \leftarrow \emph{visited} \oplus^\BbbB \emph{current\_front}$ \tcp*{Записываем вершины из фронта в посещённые}
        $\emph{new\_front} \leftarrow \emph{current\_front} \mmult{\BbbB} M$ \tcp*{Вычисляем вершины, достижимые за один шаг из текущего фронта}
        $\emph{current\_front} \leftarrow \emph{new\_front} \oplus^\BbbM \emph{visited}$ \tcp*{Формируем новый фронт с учётом уже посещённых вершин}
    }
    \KwRet{\emph{visited}}\;
    \caption{Алгоритм обхода в ширину в терминах линейной алгебры}
    \label{algo:BFS_linal}
\end{algorithm}

Заметим, что данный алгоритм решает очень простую задачу: выясняет, какие вершины достижимы из стартовой.
Решение более содержательных задач потребует, прежде всего, создания более сложных алгебраических структур.
Однако, общий принцип, шаги алгоритма, останется неизменным.

%$$\boxed{\oplus}^{\BbbB}$$
%
%{
%        \setlength{\fboxsep}{0.5\fboxsep}
%        $$\boxed{\oplus}^{\BbbB}$$
%    }
%
%    {
%        \setlength{\fboxsep}{0.4\fboxsep}
%        $$\boxed{\oplus}^{\BbbB}$$
%    }
%
%    {
%        \setlength{\fboxsep}{0.3\fboxsep}
%        $$\boxed{\oplus}^{\BbbB}$$
%    }
%
%    {
%        $$ X \Krond Y $$
%    }
%
%    {
%        \setlength{\fboxsep}{0.2\fboxsep}
%        $$X \boxed{\oplus}^{\BbbB} Y$$
%    }
%
%    {
%
%        $$X \Kronf Y$$
%    }
%
%    {
%
%        $$X \Kronw Y$$
%    }
%    {
%
%        $$X \Kronq Y$$
%    }
%
%    {
%        \setlength{\fboxsep}{0.25\fboxsep}
%        $$X \boxed{\circ}^{\BbbB} Y$$
%    }
%
%    {
%        \setlength{\fboxsep}{0.2\fboxsep}
%        $$X \boxed{\circ}^{\BbbB} Y$$
%    }
%
%$$\boxplus^{\BbbB}$$
%$$\square \oplus {\BbbB}$$
%$$\box \oplus {\BbbB}$$

\begin{example}

    Рассмотрим обход в ширину графа из примера~\ref{example:diGraph} начиная с вершины \circled{2}.
    Для обозначения текущего фронта будем использовать зелёный цвет, а для достижимых из него за один шаг~--- жёлтый, обведём вершину красным, если она посещается повторно.

    В начальном состоянии посещённых вершин нет и во фронте отмечена только вторая вершина (так как она является стартовой):
    \begin{align*}
        \emph{current\_front} & =
        \begin{pmatrix}
            0 & 0 & 1 & 0
        \end{pmatrix}
        \\
        \emph{visited}        & =
        \begin{pmatrix}
            0 & 0 & 0 & 0
        \end{pmatrix}.
    \end{align*}

    \begin{marginfigure}
        \begin{center}
            \resizebox{\marginparwidth}{!}{\input{part_01_Prep/chapter_03_GraphTheoryIntro/figures/graph_BFS_1.tex}}
        \end{center}
        \caption{Обход в ширину, шаг первый}
        \label{fig:bfs_step_1}
    \end{marginfigure}

    В посещённые вершины добавим вершины текущего фронта.
    \begin{align*}
        \emph{visited} & = \emph{visited} \oplus^\BbbB \emph{current\_front} \\
                       & =
        \begin{pmatrix}
            0 & 0 & 0 & 0
        \end{pmatrix}\oplus^\BbbB
        \begin{pmatrix}
            0 & 0 & 1 & 0
        \end{pmatrix}
    \end{align*}

    Умножим фронт на матрицу смежности для того, чтобы выяснить, в какие вершины мы можем перейти из стартовой, и получить новый фронт:
    \begin{align*}
        \emph{new\_front} & = \emph{current\_front} \mmult{\BbbB} M \\
                          & =
        \begin{pmatrix}
            0 & 0 & 1 & 0
        \end{pmatrix} \mmult{\BbbB}
        \begin{pmatrix}
            0 & 1 & 0 & 0 \\
            0 & 0 & 1 & 0 \\
            1 & 0 & 0 & 1 \\
            0 & 0 & 1 & 0
        \end{pmatrix}                                              \\
                          & =
        \begin{pmatrix}
            1 & 0 & 0 & 1
        \end{pmatrix}.
    \end{align*}

    Теперь можно создать фронт для следующей итерации с учётом нового фронта и посещённых вершин:
    \begin{align*}
        \emph{current\_front} & = \emph{new\_front} \oplus^\BbbM \emph{visited}
        \\
                              & =
        \begin{pmatrix}
            1 & 0 & 0 & 1
        \end{pmatrix}\oplus^\BbbM
        \begin{pmatrix}
            0 & 0 & 1 & 0
        \end{pmatrix}                                                          \\
                              & =
        \begin{pmatrix}
            1 & 0 & 0 & 1
        \end{pmatrix}.
    \end{align*}

    Состояние графа после первой итерации показано на рисунке~\ref{fig:bfs_step_1}: вершина \circled{2}~--- изначальное состояние фронта, а вершины \circled{1} и \circled{3} достижимы из него за один шаг.

    На втором шаге из текущего фронта мы можем попасть в вершины \circled{1} и \circled{2}.
    Сначала обновим информацию о посещённых вершинах:
    \begin{align*}
        \emph{visited} & =
        \begin{pmatrix}
            0 & 0 & 1 & 0
        \end{pmatrix}\oplus^\BbbB
        \begin{pmatrix}
            1 & 0 & 0 & 1
        \end{pmatrix} =
        \begin{pmatrix}
            1 & 0 & 1 & 1
        \end{pmatrix}.
    \end{align*}

    Затем вычислим новый фронт:
    \begin{align*}
        \emph{new\_front} & =
        \begin{pmatrix}
            1 & 0 & 0 & 1
        \end{pmatrix} \mmult{\BbbB}
        \begin{pmatrix}
            0 & 1 & 0 & 0 \\
            0 & 0 & 1 & 0 \\
            1 & 0 & 0 & 1 \\
            0 & 0 & 1 & 0
        \end{pmatrix}        \\ &=
        \begin{pmatrix}
            0 & 1 & 1 & 0
        \end{pmatrix}.
    \end{align*}

    При этом вершина \circled{2} есть в \emph{visited}, потому фронт для следующей итерации будет выглядеть следующим образом:
    \begin{align*}
        \emph{current\_front} & =
        \begin{pmatrix}
            0 & 1 & 1 & 0
        \end{pmatrix} \oplus^\BbbM
        \begin{pmatrix}
            1 & 0 & 1 & 1
        \end{pmatrix}            \\ &=
        \begin{pmatrix}
            0 & 1 & 0 & 0
        \end{pmatrix}
    \end{align*}

    \begin{marginfigure}
        \begin{center}
            \resizebox{\marginparwidth}{!}{\input{part_01_Prep/chapter_03_GraphTheoryIntro/figures/graph_BFS_2.tex}}
        \end{center}
        \caption{Обход в ширину, шаг второй}
        \label{fig:bfs_step_2}
    \end{marginfigure}

    Соответствующее состояние графа представлено на рисунке~\ref{fig:bfs_step_2}.

    Следующий, третий, шаг будет завершающим.

    \begin{align*}
        \emph{current\_front} & =
        \begin{pmatrix}
            0 & 1 & 0 & 0
        \end{pmatrix}\mmult{\BbbB}
        \begin{pmatrix}
            0 & 1 & 0 & 0 \\
            0 & 0 & 1 & 0 \\
            1 & 0 & 0 & 1 \\
            0 & 0 & 1 & 0
        \end{pmatrix}            \\ &=
        \begin{pmatrix}
            0 & 0 & 1 & 0
        \end{pmatrix}
    \end{align*}

    Из текущего фронта мы можем попасть лишь в вершину \circled{2}, а она уже была посещена.
    Соответственно, фронт для следующей итерации будет пустой и мы выйдем из цикла.

    \begin{marginfigure}
        \begin{center}
            \resizebox{\marginparwidth}{!}{\input{part_01_Prep/chapter_03_GraphTheoryIntro/figures/graph_BFS_3.tex}}
        \end{center}
        \caption{Обход в ширину, шаг третий}
        \label{fig:bfs_step_3}
    \end{marginfigure}

    Финальное состояние графа можно увидеть на рисунке~\ref{fig:bfs_step_3}.
    В итоге, результирующий вектор \emph{visited} содержит единицы на всех позициях, то есть все вершины графа достижимы из вершины \circled{2}.

\end{example}

Мы рассмотрели обход графа из одной фиксированной вершины.
Однако часто возникает необходимость совершить обход графа независимо из нескольких вершин.
При этом для каждой посещённой вершины необходимо знать, из какой именно стартовой вершины она достижима\sidenote{
    Именно поэтому мы не можем пойти наивным путём: сложить все стартовые вершины в один фронт.
}.
Данную вариацию обхода будем называть обходом в ширину с несколькими стартовыми вершинами (multiple-source BFS, MS-BFS).
Для такой постановки задачи также предложен алгоритм, основанный на линейной алгебре~\sidecite{9286186}.
Идея его такая же, как у рассмотренного выше, однако фронт~--- это уже не вектор, а матрица размера $k \times |V|$, где $k$~--- количество стартовых вершин, каждая строка которой является фронтом для одной из стартовых вершин\sidenote{
    Такая конструкция действительно является лишь естественным для алгоритмов, выраженных в терминах операций линейной алгебры, способом запустить несколько независимых обходов параллельно.
}.
Псевдокод такой вариации представлен на листинге~\ref{algo:MS-BFS_linal}.
Можно заметить, что алгоритм отличается от обхода с одной стартовой вершиной только инициализацией данных.

\begin{algorithm}
    \SetAlgoLined
    \KwData{$M$~--- булева матрица смежности графа\;
        $K$~--- вектор номеров стартовой вершины}
    \KwResult{Матрица, $i$-я строка которой указывает, какие вершины достижимы из $K[i]$}
    $\emph{current\_front} \leftarrow \Bbbzero^{k \times n}$\;
    $\emph{visited} \leftarrow \Bbbzero^{k \times n}$\;
    \For{$i \in 0\ldots |K|-1$}{$\emph{current\_front}[i,K[i]] \leftarrow 1$\;}
    \While{\emph{current\_front} $\neq \ \overline{0}$}{
        $\emph{visited} \leftarrow \emph{visited} \oplus^\BbbB \emph{current\_front}$ \tcp*{Записываем вершины из фронта в посещённые}
        $\emph{new\_front} \leftarrow \emph{current\_front} \mmult{\BbbB} M$ \tcp*{Вычисляем вершины, достижимые за один шаг из текущего фронта}
        $\emph{current\_front} \leftarrow \emph{new\_front} \oplus^\BbbM \emph{visited}$ \tcp*{Формируем новый фронт с учётом уже посещённых вершин}
    }
    \KwRet{\emph{visited}}\;
    \caption{Алгоритм поиска в ширину от нескольких стартовых вершин в терминах линейной алгебры}
    \label{algo:MS-BFS_linal}
\end{algorithm}

\begin{example}
    Рассмотрим обход в ширину графа из примера~\ref{example:diGraph} начиная с вершин \circled{1} и \circled{3}.
    Будем использовать ту же цветовую схему, что и в предыдущем примере.

    В данном случае вспомогательные структуры инициализируются следующими матрицами.
    \begin{align*}
        \emph{current\_front} & =
        \begin{pmatrix}
            0 & 1 & 0 & 0 \\
            0 & 0 & 0 & 1
        \end{pmatrix}
        \\
        \emph{visited}        & =
        \begin{pmatrix}
            0 & 0 & 0 & 0 \\
            0 & 0 & 0 & 0
        \end{pmatrix}.
    \end{align*}

    Далее будут выполняться известные уже шаги.
    Проследим изменения фронта.

    \begin{align*}
        \emph{new\_front} & = \emph{current\_front} \mmult{\BbbB} M \\
                          &=
        \begin{pmatrix}
            0 & 1 & 0 & 0 \\
            0 & 0 & 0 & 1
        \end{pmatrix} \mmult{\BbbB}
        \begin{pmatrix}
            0 & 1 & 0 & 0 \\
            0 & 0 & 1 & 0 \\
            1 & 0 & 0 & 1 \\
            0 & 0 & 1 & 0
        \end{pmatrix}        \\ &=
        \begin{pmatrix}
            0 & 0 & 1 & 0 \\
            0 & 0 & 1 & 0
        \end{pmatrix}.
    \end{align*}

    \begin{marginfigure}
        \begin{center}
            \resizebox{\marginparwidth}{!}{\input{part_01_Prep/chapter_03_GraphTheoryIntro/figures/graph_MS-BFS_1.tex}}
        \end{center}
        \caption{Обход в ширину с несколькими стартовыми вершинами, шаг первый}
        \label{fig:ms_bfs_step_1}
    \end{marginfigure}

    Первый шаг представлен на изображении~\ref{fig:ms_bfs_step_1}: из двух стартовых вершин мы можем попасть в вершину \circled{2}.
    Можно заметить, что фронт хранит информацию о том, что \circled{2} достижима и из вершины \circled{1}, и из вершины \circled{1}.

    \begin{marginfigure}
        \begin{center}
            \resizebox{\marginparwidth}{!}{\input{part_01_Prep/chapter_03_GraphTheoryIntro/figures/graph_MS-BFS_2.tex}}
        \end{center}
        \caption{Обход в ширину с несколькими стартовыми вершинами, шаг второй}
        \label{fig:ms_bfs_step_2}
    \end{marginfigure}

    Проделаем ещё один шаг.
    Он будет последим содержательным шагом: на следующем шаге мы всего лишь узнаем, что все вершины посещены.
    \begin{align*}
        \emph{new\_front} & = \emph{current\_front} \mmult{\BbbB} M \\
                          & =
        \begin{pmatrix}
            0 & 0 & 1 & 0 \\
            0 & 0 & 1 & 0
        \end{pmatrix} \mmult{\BbbB}
        \begin{pmatrix}
            0 & 1 & 0 & 0 \\
            0 & 0 & 1 & 0 \\
            1 & 0 & 0 & 1 \\
            0 & 0 & 1 & 0
        \end{pmatrix}        \\ &=
        \begin{pmatrix}
            1 & 0 & 0 & 1 \\
            1 & 0 & 0 & 1
        \end{pmatrix}.
    \end{align*}



    Результат шага представлен на рисунке~\ref{fig:ms_bfs_step_2}.
    Здесь важно обратить внимание на вершину \circled{3}: она буде повторно посещённой для стартовой вершины \circled{3}, но для вершины \circled{1} мы пришли в неё впервые.
    Потому фронт для следующей итерации будет выглядеть следующим образом.

    \begin{align*}
        \emph{current\_front} & = \emph{new\_front} \oplus^\BbbM \emph{visited} \\
                              &=
        \begin{pmatrix}
            1 & 0 & 0 & 1 \\
            1 & 0 & 0 & 1
        \end{pmatrix} \oplus^\BbbM
        \begin{pmatrix}
            0 & 1 & 1 & 0 \\
            0 & 0 & 1 & 1
        \end{pmatrix}        \\ &=
        \begin{pmatrix}
            1 & 0 & 0 & 1 \\
            1 & 0 & 0 & 0
        \end{pmatrix}.
    \end{align*}

\end{example}

Идеи, заложенные в описанных выше алгоритмах, будут использоваться нами далее, при решении других задач.

Достижимость от нескольких стартовых вершин через обход в ширину, основанный на линейной алгебре~\sidecite{9286186}
